\documentclass{standalone}
\usepackage{tikz}
\usepackage{pgfplots}

\usepgfplotslibrary{colorbrewer,colormaps}

\pgfplotsset{colormap/RdBu}
\pgfplotsset{compat=1.18}

% Match JFM fonts
\usepackage[T1]{fontenc}
\usepackage{newtxtext}
\usepackage{newtxmath}

% Final width in the JFM paper
\def\figurewidth{13.5}

% Final JFM font size
\newcommand{\fsize}{\fontsize{8.7}{10.44}\selectfont}

\begin{document}
\begin{tikzpicture}
	% \clip (\xmin,\ymin) rectangle (\xmax,\ymax);
	\def\xmin{0}
	\def\xmax{\figurewidth}
	\def\ymin{0}
	\def\ymax{1.1}

	\path[use as bounding box]
	(\xmin,\ymin) rectangle (\xmax,\ymax);

	\coordinate (clippedcenter) at ({(\xmin+\xmax)/2}, {(\ymin+\ymax)/2});

	% ============================================================
	% FIGURE-RELATIVE COLORBAR STYLE
	% ============================================================

	% Colorbar geometry
	\pgfmathsetlengthmacro{\colorbarwidth}
	{0.33*\figurewidth cm}

	\pgfmathsetlengthmacro{\colorbarheight}
	{0.08*0.33*\figurewidth cm}

	% Matplotlib-like line and tick dimensions
	\pgfmathsetlengthmacro{\colorbarlinewidth}
	{0.0020*\figurewidth cm}

	\pgfmathsetlengthmacro{\colorbarticklength}
	{0.0075*\figurewidth cm}

	\pgfplotscolorbardrawstandalone[
		colormap={example}{
				indices of colormap=(
				10,9,8,7,6,5,4,3,2,1,0 of RdBu
				),
			},
		point meta min=-0.1,
		point meta max=0.1,
		colorbar horizontal,
		samples=200,
		scaled x ticks=false,
		colorbar style={
				xtick={-0.1,-0.05,...,0.1},
				xticklabel style={
						/pgf/number format/fixed,
						/pgf/number format/precision=3,
						font=\fsize,
					},
				xlabel={$v$},
				xlabel style={
						font=\fsize,
						at={(axis description cs:0.5,-1.2)},
						anchor=north,
					},
				at={(clippedcenter)},
				anchor=south,
				yshift=0.45cm,
				width=\colorbarwidth,
				height=\colorbarheight,
				xtick pos=lower,
				xtick style={
						black,
						line width=\colorbarlinewidth,
					},
				tickwidth=\colorbarticklength,
				axis line style={
						line width=\colorbarlinewidth,
					},
			},
	]
\end{tikzpicture}
\end{document}
% 	% Colorbar outside the clipped image scope
% 	\coordinate (colorbarcenter) at (6.075,0);
%
% 	\pgfplotscolorbardrawstandalone[
% 		colormap={example}{
% 				indices of colormap=(
% 				10,9,8,7,6,5,4,3,2,1,0 of RdBu
% 				),
% 			},
% 		point meta min=-0.09,
% 		point meta max=0.09,
% 		colorbar horizontal,
% 		samples=10,
% 		scaled x ticks=false,
% 		colorbar style={
% 				xtick={-0.09,-0.045,...,0.09},
% 				xticklabel style={
% 						/pgf/number format/fixed,
% 						/pgf/number format/precision=3,
% 						font=\fsize,
% 					},
% 				xlabel={$v$},
% 				xlabel style={
% 						font=\fsize,
% 						at={(axis description cs:0.5,-1.7)},
% 						anchor=north,
% 					},
% 				at={(colorbarcenter)},
% 				yshift=0.85cm,
% 				anchor=north,
% 				width=\colorbarwidth,
% 				height=\colorbarheight,
% 				xtick pos=lower,
% 				xtick style={
% 						black,
% 						line width=\colorbarlinewidth,
% 					},
% 				tickwidth=\colorbarticklength,
% 				axis line style={
% 						line width=\colorbarlinewidth,
% 					},
% 			},
% 	]
%
% \end{tikzpicture}
%
% \end{document}
